% =====================================================================
% LAMPIRAN G: LANDASAN TEORI PELENGKAP
% =====================================================================
\chapter{Landasan Teori Pelengkap}
\label{app:teori-pelengkap}

Lampiran ini menampung uraian teoretis yang sebelumnya berada di Bab~\ref{ch:tinjauan} namun bersifat latar umum dan tidak menjadi keputusan metodologi langsung pada Bab~\ref{ch:metode}. Pemindahan dilakukan agar Bab~\ref{ch:tinjauan} tetap padat dan terfokus pada teori yang dipakai pada \textit{pipeline} akhir, yaitu \textit{encoder} DINOv3, \textit{multilingual-E5}, fusi awal, dan \textit{classifier head} CatBoost.

% =====================================================================
\section*{G.1 \textit{Web Scraping}: Klasifikasi Halaman, Alur, dan Pustaka}
\label{app:webscraping-detail}

Halaman web yang menjadi target \textit{scraping} dibedakan menjadi dua kelompok berdasarkan cara konten dirender. \textbf{Halaman statis} adalah halaman yang seluruh kontennya sudah terbentuk pada respons HTML server, sehingga cukup memerlukan pengunduh HTTP dan pengurai HTML. \textbf{Halaman dinamis} adalah halaman yang kontennya dirender oleh JavaScript pada peramban (\textit{client-side rendering}), sehingga memerlukan eksekusi JavaScript melalui peramban tanpa kepala (\textit{headless browser}) atau pemanggilan API XHR yang ditemukan melalui inspeksi jaringan.

Alur umum \textit{web scraping} terdiri atas lima tahap, yaitu (1) identifikasi sumber (halaman daftar, halaman rincian, dan pola URL), (2) pengunduhan melalui permintaan HTTP GET dengan memperhatikan \textit{user-agent}, \textit{cookie}, dan pembatasan kecepatan, (3) penguraian (\textit{parsing}) elemen HTML yang relevan, (4) pembersihan dan normalisasi karakter, tanggal, serta pemetaan label bebas ke kategori standar, dan (5) penyimpanan data terstruktur (CSV, JSONL, Parquet) atau basis data relasional.

Beberapa pustaka populer pada ekosistem Python untuk \textit{web scraping} antara lain \texttt{requests} (pustaka HTTP sederhana untuk halaman statis), \texttt{BeautifulSoup} (pengurai HTML yang fleksibel), \texttt{lxml} (pengurai cepat berbasis libxml2 dengan dukungan XPath), \texttt{Scrapy} (\textit{framework} berbasis \textit{spider} dan \textit{pipeline} untuk \textit{scraping} berskala besar), \texttt{Selenium} dan \texttt{Playwright} (otomatisasi peramban yang dapat mengeksekusi JavaScript pada halaman dinamis), serta \texttt{httpx} (alternatif modern \texttt{requests} dengan dukungan asinkron).

% =====================================================================
\section*{G.2 \textit{Multi-Head Attention}: Formulasi Lengkap}
\label{app:mha-detail}

\textit{Multi-head attention} adalah perluasan \textit{single-head attention} yang menjalankan beberapa proyeksi linear paralel dari \textit{query}, \textit{key}, dan \textit{value}, kemudian menggabungkan keluarannya. Motivasi desain ini adalah bahwa satu \textit{head} terbatas pada satu pola hubungan antar token, sementara teks dan citra alami memuat banyak pola hubungan secara simultan.

Misalkan dimensi \textit{embedding} model adalah $d_{\text{model}}$ dan jumlah \textit{head} adalah $h$. Setiap \textit{head} bekerja pada subruang berdimensi $d_k = d_{\text{model}} / h$. Untuk setiap \textit{head} $i \in \{1, \dots, h\}$ didefinisikan tiga matriks proyeksi yang dapat dilatih:

\begin{equation}
    \mathbf{Q}_i = \mathbf{Q} \mathbf{W}_i^Q, \quad
    \mathbf{K}_i = \mathbf{K} \mathbf{W}_i^K, \quad
    \mathbf{V}_i = \mathbf{V} \mathbf{W}_i^V,
    \label{eq:mha_proj}
\end{equation}
dengan $\mathbf{W}_i^Q, \mathbf{W}_i^K \in \mathbb{R}^{d_{\text{model}} \times d_k}$ dan $\mathbf{W}_i^V \in \mathbb{R}^{d_{\text{model}} \times d_v}$. Setiap \textit{head} menjalankan \textit{scaled dot-product attention} secara independen:

\begin{equation}
    \text{head}_i = \text{softmax}\!\left( \frac{\mathbf{Q}_i \mathbf{K}_i^\top}{\sqrt{d_k}} \right) \mathbf{V}_i.
    \label{eq:mha_head}
\end{equation}

Keluaran $h$ \textit{head} dikonkatenasi sepanjang dimensi terakhir, kemudian diproyeksikan kembali ke ruang berdimensi $d_{\text{model}}$ melalui matriks proyeksi keluaran $\mathbf{W}^O \in \mathbb{R}^{h d_v \times d_{\text{model}}}$:

\begin{equation}
    \text{MultiHead}(\mathbf{Q}, \mathbf{K}, \mathbf{V}) = \text{Concat}(\text{head}_1, \dots, \text{head}_h) \, \mathbf{W}^O.
    \label{eq:multihead}
\end{equation}

Dimensi $d_k$ pada setiap \textit{head} biasanya jauh lebih kecil daripada $d_{\text{model}}$, contohnya $d_{\text{model}} = 768$ dengan $h = 12$ menghasilkan $d_k = 64$, sehingga total komputasi tetap sebanding dengan \textit{single-head attention} berdimensi $d_{\text{model}}$. Pada model bahasa modern jumlah \textit{head} umumnya berkisar antara 8 dan 32, dan penambahan \textit{head} melampaui jumlah optimal sering kali tidak meningkatkan performa serta membuat sebagian \textit{head} menjadi redundan.

% =====================================================================
\section*{G.3 \textit{Backbone} Visual Pembanding: EVA-02 dan Hiera}
\label{app:backbone-detail}

\subsection*{G.3.1 EVA-02}

EVA-02 \citep{fangEVA022023} adalah \textit{vision foundation model} generasi kedua dari seri EVA yang menggabungkan \textit{Masked Image Modeling} (MIM) dan penyelarasan dengan representasi CLIP. Inovasi utamanya adalah penggunaan fitur CLIP terlatih sebagai target sinyal supervisi dalam proses MIM, alih-alih piksel. Pada sisi arsitektur, EVA-02 memperbaiki ViT melalui Sub-LN (penambahan \textit{layer normalization} di dalam blok \textit{attention} dan FFN), SwiGLU FFN (varian \textit{gated} pengganti FFN GELU), dan 2D RoPE (\textit{rotary position embedding} dua dimensi). EVA-02 dilatih pada dataset Merged-2B dengan target fitur dari EVA-CLIP, dan menempati posisi kompetitif terhadap \textit{state-of-the-art} pada \textit{benchmark} ImageNet, COCO, dan ADE20K.

\subsection*{G.3.2 Hiera}

Hiera (\textit{Hierarchical Vision Transformer without the Bells and Whistles}) \citep{ryaliHiera2023} menyederhanakan ViT dengan menghapus komponen yang dianggap tidak diperlukan, berlandaskan hipotesis bahwa banyak komponen pada \textit{hierarchical Transformer} sebelumnya (Swin, MViT) menjadi tidak perlu apabila pelatihan dilakukan dengan paradigma MAE. Inovasi utamanya adalah penghapusan \textit{positional embedding} eksplisit dan penggunaan \textit{pooling attention} hierarkis yang terinspirasi MAE \citep{heMAE2022}, sehingga menghasilkan representasi hierarkis menyerupai piramida fitur CNN namun tetap dalam kerangka \textit{Transformer} murni. \textit{Masked Autoencoder} (MAE) sendiri menutup sekitar 75\% \textit{patch} citra dan meminta model merekonstruksi piksel pada \textit{patch} yang ditutup; asimetri antara \textit{encoder} (melihat \textit{patch} terlihat) dan \textit{decoder} (merekonstruksi semua \textit{patch}) memberikan efisiensi komputasi yang signifikan.

% =====================================================================
\section*{G.4 \textit{Encoder} Teks Pembanding}
\label{app:textencoder-detail}

\subsection*{G.4.1 BERT dan mBERT}

BERT \citep{devlinBERT2019} adalah \textit{Transformer encoder} dua arah yang dilatih dengan \textit{Masked Language Modeling} (MLM) dan \textit{Next Sentence Prediction} (NSP), menggunakan tokenizer WordPiece dengan kosakata sekitar 30 ribu pada varian Inggris. \textit{Multilingual BERT} (mBERT) adalah varian yang dilatih pada Wikipedia 104 bahasa secara serentak dengan kosakata WordPiece bersama berukuran sekitar 110 ribu, sehingga memperoleh kemampuan transfer lintas bahasa \textit{zero-shot}.

\subsection*{G.4.2 IndoBERT}

IndoBERT \citep{kotoIndoLEM2020,wilieIndoNLU2020} adalah keluarga \textit{Transformer encoder} yang dilatih khusus pada korpus Bahasa Indonesia melalui dua jalur publikasi, yaitu IndoLEM/IndoBERT \citep{kotoIndoLEM2020} dan IndoNLU/IndoBERT \citep{wilieIndoNLU2020}. Arsitekturnya identik dengan BERT, namun dengan tokenizer dan korpus pra-latih spesifik Bahasa Indonesia, sehingga menghasilkan representasi yang lebih kaya pada idiom, ejaan, dan istilah lokal dibanding mBERT.

\subsection*{G.4.3 Cendol-mT5}

Cendol \citep{cahyawijayaCendol2024} adalah keluarga model bahasa berbasis instruksi yang dirilis oleh IndoNLP untuk Bahasa Indonesia dan bahasa daerah Nusantara. Sub-keluarga Cendol-mT5 dibangun di atas arsitektur mT5 (\textit{encoder-decoder}); pada penggunaan sebagai \textit{encoder} kalimat, hanya bagian \textit{encoder} mT5 yang dimanfaatkan dengan \textit{mean pooling} yang dinormalisasi. mT5 \citep{xueMT52021} adalah varian multibahasa T5 yang dilatih pada korpus mC4 mencakup 101 bahasa dengan tujuan \textit{span corruption}.

\subsection*{G.4.4 BGE-M3}

BGE-M3 \citep{chenBGEM32024} adalah \textit{embedding model} multibahasa dari BAAI yang mendukung multibahasa (100+ bahasa), multifungsi (\textit{dense}, \textit{sparse}, dan \textit{multi-vector retrieval}), serta multigranularitas (hingga 8.192 token). Arsitekturnya didasarkan pada XLM-RoBERTa \textit{large} (sekitar 560 juta parameter) dengan strategi \textit{self-knowledge distillation}. \textit{Embedding} yang dihasilkan berdimensi 1024 dari token CLS dengan normalisasi L2.

% =====================================================================
\section*{G.5 \textit{Unified Modeling Language} (UML): Jenis Diagram dan C4 Model}
\label{app:uml-detail}

UML \citep{boochUML2005} adalah bahasa permodelan visual standar yang dibakukan oleh \textit{Object Management Group} (OMG). Versi UML 2.x mendefinisikan empat belas jenis diagram yang dikelompokkan menjadi diagram struktur dan diagram perilaku.

\textbf{Diagram struktur} memodelkan elemen sistem yang relatif statis. Tiga jenis yang paling sering dipakai adalah Diagram Kelas (\textit{class diagram}: kelas, atribut, metode, dan relasi), Diagram Komponen (\textit{component diagram}: komponen yang dapat di-\textit{deploy} beserta antarmukanya), dan Diagram Deployment (\textit{deployment diagram}: pemetaan komponen ke simpul perangkat keras).

\textbf{Diagram perilaku} memodelkan dinamika sistem. Empat jenis yang relevan adalah Diagram \textit{Use Case} (interaksi aktor dengan fungsionalitas), Diagram \textit{Sequence} (urutan pesan antar objek pada sumbu waktu), Diagram Aktivitas (\textit{activity diagram}: alur kerja dengan simpul keputusan dan paralel), dan Diagram \textit{State Machine} (keadaan objek dan transisinya).

Sebagai pelengkap, \citet{brownC42018} memperkenalkan \textit{C4 Model} yang menyederhanakan dokumentasi arsitektur menjadi empat tingkat abstraksi, yaitu \textit{Context}, \textit{Container}, \textit{Component}, dan \textit{Code}. C4 sering dipadukan dengan UML dan lebih mudah dibaca oleh pemangku kepentingan non-teknis.

% =====================================================================
\section*{G.6 Docker dan Kontainerisasi: Konsep Rinci}
\label{app:docker-detail}

Docker \citep{merkelDocker2014} adalah \textit{platform} kontainerisasi yang mengemas aplikasi beserta dependensinya ke dalam satu unit yang konsisten di berbagai lingkungan. Berbeda dari mesin virtual, kontainer berbagi \textit{kernel} sistem operasi induk dan mengisolasi proses melalui \textit{Linux namespaces} dan \textit{cgroups}, sehingga ringan dan cepat dijalankan.

Lima konsep inti Docker adalah \textbf{Image} (\textit{template} \textit{read-only} berisi sistem berkas aplikasi), \textbf{Container} (instansi \textit{image} yang berjalan), \textbf{Dockerfile} (skrip deklaratif langkah \textit{build}), \textbf{Volume} (penyimpanan persisten), dan \textbf{Registry} (gudang \textit{image} terpusat seperti Docker Hub). Sebuah \texttt{Dockerfile} umumnya memuat instruksi \texttt{FROM}, \texttt{WORKDIR}, \texttt{COPY}/\texttt{ADD}, \texttt{RUN}, \texttt{EXPOSE}, \texttt{ENV}, \texttt{HEALTHCHECK}, serta \texttt{CMD}/\texttt{ENTRYPOINT}.

Strategi \textit{multi-stage build} memisahkan tahap \textit{build} (yang membutuhkan kompilator) dari tahap \textit{runtime} (yang hanya membutuhkan binari hasil), sehingga \textit{image} akhir lebih ramping. Docker Compose mendefinisikan aplikasi multi-kontainer melalui berkas YAML, sehingga lingkungan (server inferensi, basis data, \textit{reverse proxy}) dapat dijalankan dengan satu perintah \texttt{docker compose up}. Pada \textit{pipeline machine learning}, Docker memberikan keuntungan berupa \textit{reproducibility}, isolasi dependensi (PyTorch, ONNX Runtime, CUDA), \textit{portability}, dan akselerasi GPU melalui \texttt{nvidia-container-toolkit}.
